%% annex/p03/types-ultraproducts.tex --- proof annex fragment for
%% parts/p03/types-ultraproducts.tex.  Source-complete proofs are compressed
%% individually toward the Phase-4 40--60 percent target while retaining every
%% load-bearing step, including repairs R-028, R-029, and R-036.

\begin{annexitem}{annex:p3.types.1}{thm:p03-type-space-compact}{Compactness of type spaces}
\begin{proof}
The displayed basis is clopen.  If \(p\ne q\), choose
\(\varphi\in p\setminus q\); completeness puts \(\neg\varphi\in q\), so
\([\varphi]\) and \([\neg\varphi]\) separate the points.

It remains to treat a basic open cover \(\{[\varphi_i]:i\in I\}\), since
every open cover has a basic refinement.  If no finite subfamily covers, then
\[
 \Sigma=T\cup\{\neg\varphi_i(\bar d):i\in I\}
\]
is finitely satisfiable: for finite \(J\subseteq I\), choose
\(p\in S_n(T)\setminus\bigcup_{j\in J}[\varphi_j]\); a realization of
\(p\) in some model of \(T\) satisfies the corresponding finite subset of
\(\Sigma\).  First-order compactness gives a model \(\mathcal N\) of
\(\Sigma\), with \(\bar b\) interpreting \(\bar d\).

Now perform the complete-type step required by R-028: set
\(q=\{\psi(\bar x):\mathcal N\entails\psi(\bar b)\}\).  Exactly one of
\(\psi,\neg\psi\) belongs to \(q\), and \((\mathcal N,\bar b)\) witnesses
its consistency with \(T\); hence \(q\in S_n(T)\).  But
\(\mathcal N\entails\neg\varphi_i(\bar b)\) for every \(i\), so \(q\) lies
in no member of the cover, a contradiction.  Thus every basic, and hence
every open, cover has a finite subcover.
\end{proof}
\end{annexitem}

\begin{annexitem}{annex:p3.types.2}{thm:p03-saturation-homogeneity}{Saturation supplies homogeneity}
\begin{proof}
Let \(h:A\to\dom\mathcal M\) be elementary, \(\abs A<\kappa\), and
\(b\in\dom\mathcal M\).  If \(b\in A\), use \(h\).  Otherwise let
\(p(x)\) be the complete type of \(b\) over \(A\), and let \(q(x)=h(p)\)
be obtained by replacing every parameter \(a\) by \(h(a)\).

The consistency of \(q\) is the compactness step required by R-029.  Any
finite part has the form \(h(\varphi_1),\ldots,h(\varphi_m)\) with
\(\varphi_j\in p\).  Since \(b\) witnesses
\(\exists x\bigwedge_j\varphi_j(x)\) and \(h\) is elementary,
\(\mathcal M\) satisfies
\(\exists x\bigwedge_jh(\varphi_j)(x)\).  Interpreting one new constant by
such a witness satisfies that finite part together with the complete
\(\Lang(h[A])\)-theory of \(\mathcal M_{h[A]}\).  First-order compactness
therefore makes \(q\)
consistent with \(\mathcal M\).  Parameter replacement commutes with
negation and reaches every formula over \(h[A]\), so \(q\) is complete.

Because \(\abs{h[A]}\leq\abs A<\kappa\), saturation gives
\(d\in\dom\mathcal M\) realizing \(q\).  Extend \(h\) by \(h'(b)=d\).
For each formula \(\chi(x,\bar a)\) over \(A\), completeness of \(p\) and
realization of \(q\) give
\[
 \mathcal M\entails\chi(b,\bar a)
 \quad\Longleftrightarrow\quad
 \mathcal M\entails\chi(d,h(\bar a)).
\]
Replacing repeated occurrences of \(b\) by one variable reduces every tuple
from \(A\cup\{b\}\) to this equivalence; hence \(h'\) is elementary.
\end{proof}
\end{annexitem}

\begin{annexitem}{annex:p3.types.3}{thm:p03-nonprincipal-ultrafilter}{Existence and characterization of nonprincipal ultrafilters}
\begin{proof}
The cofinite family contains \(I\), omits \(\varnothing\), is closed under
intersection because a union of two finite complements is finite, and is
upward closed; it is therefore the Fr\'echet filter.  The Ultrafilter Lemma
extends it to an ultrafilter \(\mathcal U\).  Such an extension cannot be
principal at \(i\), because it contains the cofinite set
\(I\setminus\{i\}\).

It remains to prove the converse characterization, including the finite-set
argument required by R-036.  If \(X\cup Y\in\mathcal U\), then at least one
of \(X,Y\) lies in \(\mathcal U\): otherwise both complements would lie in
\(\mathcal U\), forcing the complement of \(X\cup Y\) into the filter.
Induction on the size of a nonempty finite \(F\in\mathcal U\) therefore puts
some singleton \(\{i\}\) in \(\mathcal U\).  Upward closure and exclusion
of the empty intersection then show
\(\mathcal U=\{X\subseteq I:i\in X\}\), so \(\mathcal U\) is principal.
Contrapositively, a nonprincipal ultrafilter contains no finite set; the
ultrafilter alternative consequently puts the complement of every finite set
in it.  Conversely, an ultrafilter containing all cofinite sets extends the
Fr\'echet filter and, by the first paragraph, is nonprincipal.  Taking
\(I=\N\) yields the final claim.
\end{proof}
\end{annexitem}

\begin{annexitem}{annex:p3.types.4}{thm:p03-ultraproduct-well-defined}{Ultraproduct well-definedness}
\begin{proof}
Reflexivity, symmetry, and constants are immediate;
transitivity uses intersection of large equality sets.  Equivalent
representatives agree on large intersection, so function values agree there
and relation truth sets have identical ultrafilter membership.
\end{proof}
\end{annexitem}

\begin{annexitem}{annex:p3.types.5}{thm:p03-los-terms}{Term evaluation in an ultraproduct}
\begin{proof}
Structural induction: variables and constants are coordinatewise by
definition.  For \(t=f(t_1,\ldots,t_m)\), apply the hypotheses to the immediate
subterms and then the coordinatewise interpretation of \(f\).
\end{proof}
\end{annexitem}

\begin{annexitem}{annex:p3.types.6}{thm:p03-los}{\L o\'s's theorem}
\begin{proof}
For a formula \(\theta(\bar x)\) and representatives \(\bar a\), write
\[
 X_\theta(\bar a)=
 \{i\in I:\mathcal M_i\entails\theta(\bar a(i))\}.
\]
We prove the equivalence by structural induction.  For \(t=s\), term
evaluation and the quotient definition say that the ultraproduct satisfies
the equality exactly when \(X_{t=s}\in\mathcal U\).  Relation atoms follow
from the same term lemma and the definition of ultraproduct relations;
\(X_\bot=\varnothing\notin\mathcal U\).

For the Boolean cases,
\(X_{\neg\alpha}=I\setminus X_\alpha\) and
\(X_{\alpha\land\beta}=X_\alpha\cap X_\beta\).  The ultrafilter alternative
and filter closure therefore match the truth clauses for negation and
conjunction.  Also
\(X_{\alpha\lor\beta}=X_\alpha\cup X_\beta\); an ultrafilter contains this
union exactly when it contains a member, since otherwise it contains both
complements.  Finally
\(X_{\alpha\to\beta}=(I\setminus X_\alpha)\cup X_\beta\), which is large
exactly when \(X_\alpha\) is not large or \(X_\beta\) is large.  The
induction hypotheses finish every Boolean connective.

Let \(\theta(\bar x)=\exists y\,\psi(y,\bar x)\).  If the ultraproduct has a
witness \([b]_{\mathcal U}\), the induction hypothesis makes
\(X_\psi(b,\bar a)\) large; it is contained in \(X_\theta(\bar a)\), so the
latter is large.  Conversely, if \(X_\theta(\bar a)\) is large, choose in
each of its coordinates a witness \(b(i)\), and choose an arbitrary element
of each nonempty remaining domain elsewhere.  Choice supplies this function.
Then \(X_\psi(b,\bar a)\supseteq X_\theta(\bar a)\) is large, and the
induction hypothesis makes \([b]_{\mathcal U}\) an ultraproduct witness.
Universal formulas follow from the already proved negation and existential
cases using \(\forall y\,\psi\) iff \(\neg\exists y\,\neg\psi\).  This
completes the induction.
\end{proof}
\end{annexitem}

\begin{annexitem}{annex:p3.types.7}{thm:p03-diagonal-ultrapower}{Diagonal elementary embedding}
\begin{proof}
Constant truth sets are \(I\) or \(\varnothing\);
Theorem~\ref{thm:p03-los} makes \(d\) elementary.
\end{proof}
\end{annexitem}

\begin{annexitem}{annex:p3.types.8}{thm:p03-compactness-ultraproduct}{Compactness via ultraproducts}
\begin{proof}
Let \(I\) be the set of finite subsets of \(\Gamma\).  By Choice, select for
each \(\Delta\in I\) a \(\Lang\)-structure
\(\mathcal M_\Delta\entails\Delta\).  For \(\sigma\in\Gamma\), put
\[
 I_\sigma=\{\Delta\in I:\sigma\in\Delta\}.
\]
The family \(\{I_\sigma:\sigma\in\Gamma\}\) has the finite-intersection
property: for \(\sigma_1,\ldots,\sigma_m\), the index
\(\{\sigma_1,\ldots,\sigma_m\}\) belongs to every corresponding
\(I_{\sigma_j}\).  Hence its upward, finite-intersection closure is a proper
filter.  The Ultrafilter Lemma extends it to an ultrafilter \(\mathcal U\).

For each \(\sigma\in\Gamma\), the generated filter and its extension give
\(I_\sigma\in\mathcal U\).  Every \(\Delta\in I_\sigma\) satisfies \(\sigma\),
so upward closure gives
\[
 I_\sigma\subseteq
 \{\Delta\in I:\mathcal M_\Delta\entails\sigma\},
 \qquad
 \{\Delta\in I:\mathcal M_\Delta\entails\sigma\}\in\mathcal U.
\]
By \L o\'s's theorem,
\(\prod_{\Delta\in I}\mathcal M_\Delta/\mathcal U\entails\sigma\).
This holds for every \(\sigma\in\Gamma\), so the ultraproduct is a model of
\(\Gamma\).
\end{proof}
\end{annexitem}
